\documentclass[11pt,a4paper]{article}
\usepackage{fontspec}
\IfFontExistsTF{Times New Roman}{\setmainfont{Times New Roman}}{\setmainfont{Tinos}}
\IfFontExistsTF{Consolas}{\setmonofont{Consolas}[Scale=0.90]}{\setmonofont{Liberation Mono}[Scale=0.90]}
\usepackage[english]{babel}
\usepackage[top=1.35cm,bottom=1.35cm,left=1.65cm,right=1.65cm]{geometry}
\usepackage{xcolor,enumitem,listings,booktabs,array,amsmath,amssymb,titlesec,fancyhdr,parskip,tabularx,lastpage,hyperref}
\usepackage{tcolorbox}
\tcbuselibrary{skins}

\definecolor{primaryblue}{RGB}{18,65,115}
\definecolor{secondaryteal}{RGB}{14,102,110}
\definecolor{CodeBg}{RGB}{248,250,252}
\definecolor{DarkGray}{RGB}{55,55,55}
\definecolor{boxbg}{RGB}{248,250,252}

\setlength{\parindent}{0pt}\setlength{\parskip}{0pt}
\setlist[itemize]{leftmargin=1.3em,itemsep=1pt,topsep=1pt,parsep=0pt,partopsep=0pt}
\setlist[enumerate]{leftmargin=1.4em,itemsep=1pt,topsep=1pt,parsep=0pt,partopsep=0pt}
\setlist[enumerate,2]{topsep=0.5pt,itemsep=0pt,parsep=0pt,partopsep=0pt}

\lstset{
  language=Python,
  basicstyle=\ttfamily\scriptsize,
  keywordstyle=\color{blue!75!black}\bfseries,
  stringstyle=\color{red!65!black},
  commentstyle=\color{primaryblue!80!black}\itshape,
  showstringspaces=false,
  columns=fullflexible,
  keepspaces=true,
  breaklines=true,
  frame=single,
  rulecolor=\color{primaryblue!45},
  backgroundcolor=\color{CodeBg},
  xleftmargin=0.2em,
  xrightmargin=0.2em,
  aboveskip=0.15em,
  belowskip=0.15em,
  tabsize=4
}

\pagestyle{fancy}\fancyhf{}\setlength{\headheight}{14pt}
\fancyhead[L]{\footnotesize\textbf{DSAI1003 -- Fundamental Programming Concepts in Python}}
\fancyhead[R]{\footnotesize\textbf{Inheritance \& Polymorphism Exam -- 60 Minutes}}
\fancyfoot[L]{\scriptsize Faculty of Data Science and Artificial Intelligence -- National Economics University}
\fancyfoot[R]{\footnotesize Page \thepage\ of \pageref{LastPage}}
\renewcommand{\headrulewidth}{0.4pt}\renewcommand{\footrulewidth}{0.3pt}

\titleformat{\section}{\large\bfseries\color{primaryblue}}{}{0pt}{}
\titleformat{\subsection}{\normalsize\bfseries\color{primaryblue}}{}{0pt}{}
\titlespacing*{\section}{0pt}{2.5pt}{1pt}\titlespacing*{\subsection}{0pt}{1.5pt}{1pt}

\newtcolorbox{headerbox}{colback=white,colframe=primaryblue,boxrule=0.7pt,arc=1.5mm,left=2.5mm,right=2.5mm,top=1mm,bottom=1mm}
\newtcolorbox{answerbox}[2][]{enhanced,colback=white,colframe=DarkGray!45,boxrule=0.5pt,arc=1mm,left=2mm,right=2mm,top=1mm,bottom=1mm,title=\footnotesize\color{DarkGray}\textbf{#2},#1}

\hypersetup{
  colorlinks=true,
  linkcolor=primaryblue,
  urlcolor=primaryblue,
  pdftitle={DSAI1003 - Inheritance and Polymorphism Progress Exam},
  pdfauthor={Faculty of Data Science and Artificial Intelligence - NEU}
}

\begin{document}
% =========================================================================
% PAGE 1: INSTITUTIONAL HEADER, ANSWER SHEET, SECTION A (Q1-Q5)
% =========================================================================
\noindent
\begin{minipage}[t]{0.63\textwidth}\raggedright\small
\textbf{NATIONAL ECONOMICS UNIVERSITY}\\
\textbf{COLLEGE OF TECHNOLOGY}\\
\textbf{FACULTY OF DATA SCIENCE \& ARTIFICIAL INTELLIGENCE (FDA)}
\end{minipage}\hfill
\begin{minipage}[t]{0.35\textwidth}\raggedleft\footnotesize
\textbf{PROGRESS EXAMINATION}\\
Academic Year: 2026--2027\\
Semester: Fall 2026 \quad \textbf{DSAI1003}
\end{minipage}
\vspace{0.1em}
\begin{center}
{\Large\bfseries\color{primaryblue} FUNDAMENTAL PROGRAMMING CONCEPTS IN PYTHON}\\[0.05em]
{\normalsize\bfseries PROGRESS EXAM: INHERITANCE AND POLYMORPHISM (CHAPTER 10)}\\[0.05em]
{\normalsize\textbf{Paper Code:} \rule{1.5cm}{0.4pt} \quad | \quad \textbf{Duration: 60 Minutes} (Excluding distribution)}
\end{center}
\begin{headerbox}
\begin{tabularx}{\linewidth}{@{}X X@{}}
\textbf{Candidate Full Name:} \dotfill & \textbf{Student ID (MSSV):} \dotfill \\[0.15em]
\textbf{Class / Cohort:} \dotfill & \textbf{Exam Room / Seat No.:} \dotfill
\end{tabularx}
\vspace{1.5mm}
\centering\setlength{\tabcolsep}{4pt}\renewcommand{\arraystretch}{1.05}
\begin{tabular}{|l|c|c|c|c||c|c|}\hline
\textbf{Section} & \textbf{A} & \textbf{B} & \textbf{C} & \textbf{D} & \textbf{Total (100)} & \textbf{Examiner Signature} \\ \hline
\textbf{Max Score} & 20 & 25 & 25 & 30 & \textbf{100} & \\[0.2em] \cline{1-6}
\textbf{Score} & & & & & & \\[0.35em] \hline
\end{tabular}
\vspace{1mm}\raggedright\scriptsize
\textbf{Instructions:} Closed-book examination; no electronic devices or Internet access permitted. Answer all questions. Unless otherwise specified, assume Python 3 standard runtime semantics. Standard built-in libraries (such as \texttt{math}) are permitted; third-party packages (such as NumPy, Pandas) are strictly prohibited. For program-design tasks, implement modular class hierarchies with clean encapsulation, explicit superclass delegation, and a \texttt{main()} demonstration driver. Do not detach pages.
\end{headerbox}
\vspace{1mm}
\textbf{MULTIPLE-CHOICE ANSWER SHEET} \textit{(Write A, B, C, or D in each cell.)}
\begin{center}\setlength{\tabcolsep}{5pt}\renewcommand{\arraystretch}{1.1}
\begin{tabular}{|c|c|c|c|c|c|c|c|c|c|c|}\hline
\textbf{Question} & \textbf{1} & \textbf{2} & \textbf{3} & \textbf{4} & \textbf{5} & \textbf{6} & \textbf{7} & \textbf{8} & \textbf{9} & \textbf{10} \\ \hline
\textbf{Answer} & & & & & & & & & & \\[1.5mm] \hline
\end{tabular}\end{center}
\vspace{-2.5mm}
\section*{Section A: Inheritance \& Polymorphism Foundations \hfill \normalsize\normalfont(20 points)}
\textit{Answer all 10 questions. Each correct answer is worth 2 points. Suggested time: 12 minutes.}
\begin{enumerate}[label=\textbf{Q\arabic*.},leftmargin=2em,itemsep=1.5pt,topsep=1pt]
\item In object-oriented programming, which relationship correctly exemplifies \textbf{inheritance} (the ``is-a'' relationship)?
\begin{enumerate}[label=\Alph*.,leftmargin=2em]
\item A \texttt{Car} maintains an internal attribute referencing an instance of \texttt{Engine}.
\item A \texttt{Car} holds a list containing four \texttt{Tire} objects.
\item A \texttt{Car} is a specialized subclass inheriting state and behaviors from base class \texttt{Vehicle}.
\item A \texttt{Car} stores a dictionary mapping service dates to maintenance invoices.
\end{enumerate}
\item What is the standard Python syntax for declaring a subclass \texttt{SavingsAccount} that inherits from \texttt{BankAccount}?
\begin{enumerate}[label=\Alph*.,leftmargin=2em]
\item \texttt{class SavingsAccount extends BankAccount:} \quad \Alph*.~ \texttt{class SavingsAccount inherits BankAccount:}
\item \texttt{class BankAccount(SavingsAccount):} \quad \Alph*.~ \texttt{class SavingsAccount(BankAccount):}
\end{enumerate}
\item When a subclass defines its own constructor \texttt{\_\_init\_\_}, why is it essential to invoke \texttt{super().\_\_init\_\_(...)}?
\begin{enumerate}[label=\Alph*.,leftmargin=2em]
\item It allocates dynamic heap memory specifically reserved for base class bytecode.
\item It executes the superclass constructor to properly initialize inherited instance attributes.
\item The Python interpreter automatically deletes base methods if \texttt{super().\_\_init\_\_()} is omitted.
\item It prevents any external client code from instantiating the superclass.
\end{enumerate}
\item What occurs when a subclass defines a method with the \textbf{identical name and signature} as a method in its superclass?
\begin{enumerate}[label=\Alph*.,leftmargin=2em]
\item The subclass method overrides (replaces) the superclass implementation for instances of the subclass.
\item A \texttt{NameError} is raised at runtime due to duplicate method declarations in the hierarchy.
\item The Python runtime executes both the superclass method and subclass method concurrently in threads.
\item The superclass method is permanently erased from memory for instances of all classes.
\end{enumerate}
\item In a \texttt{Manager} subclass that overrides \texttt{get\_salary(self)}, why must it call \texttt{super().get\_salary()} rather than \texttt{self.get\_salary()}?
\begin{enumerate}[label=\Alph*.,leftmargin=2em]
\item Calling \texttt{self.get\_salary()} causes infinite recursive calls to \texttt{Manager.get\_salary()}, raising \texttt{RecursionError}.
\item Calling \texttt{super().get\_salary()} automatically converts the monetary output into a float.
\item Python syntax strictly forbids the keyword \texttt{self} from being referenced inside overridden methods.
\item Calling \texttt{self.get\_salary()} forces the calculation to run in restricted kernel sandbox mode.
\end{enumerate}
\end{enumerate}
\newpage

% =========================================================================
% PAGE 2: Q6-Q10 AND SECTION B Q11
% =========================================================================
\begin{enumerate}[label=\textbf{Q\arabic*.},leftmargin=2em,itemsep=1.5pt,topsep=1pt,resume]
\item What is the core capability provided by \textbf{polymorphism} in an object-oriented Python program?
\begin{enumerate}[label=\Alph*.,leftmargin=2em]
\item It forces all subclasses to store identical instance attribute values across memory addresses.
\item It restricts method invocations so that only base class definitions can ever be executed at runtime.
\item It converts diverse objects into serialized JSON strings prior to function processing.
\item It allows code to treat objects of different subclasses uniformly, dynamically invoking overridden methods at runtime based on the actual object type.
\end{enumerate}
\item Given \texttt{class Employee: pass}, \texttt{class Manager(Employee): pass}, and \texttt{class Executive(Manager): pass}. If \texttt{boss = Executive()}, which expression evaluates to \textbf{False}?
\begin{enumerate}[label=\Alph*.,leftmargin=2em]
\item \texttt{isinstance(boss, Executive)} \quad \Alph*.~ \texttt{isinstance(boss, Employee)} \quad \Alph*.~ \texttt{type(boss) is Employee} \quad \Alph*.~ \texttt{isinstance(boss, Manager)}
\end{enumerate}
\item Given base class \texttt{Question} and derived class \texttt{ChoiceQuestion(Question)}, which expression evaluates to \textbf{True}?
\begin{enumerate}[label=\Alph*.,leftmargin=2em]
\item \texttt{issubclass(Question, ChoiceQuestion)} \quad \Alph*.~ \texttt{issubclass(ChoiceQuestion, Question)} \quad \Alph*.~ \texttt{issubclass(Question, object) is False}
\end{enumerate}
\item In standard Python class design, how does a base class author conventionally enforce that subclasses must implement a template method?
\begin{enumerate}[label=\Alph*.,leftmargin=2em]
\item Define the method with body: \lstinline|raise NotImplementedError("Subclasses must override this method")|.
\item Prepend the method header with the mandatory keyword: \texttt{abstract def method(self):}.
\item Enclose the method name in quadruple underscores: \texttt{def \_\_\_\_process\_\_\_\_(self):}.
\item Omit the method entirely from the superclass so Python auto-generates it at runtime.
\end{enumerate}
\item According to Horstmann Section 10.1, which scenario represents an \textbf{inappropriate} use of inheritance?
\begin{enumerate}[label=\Alph*.,leftmargin=2em]
\item Defining \texttt{SavingsAccount} as a subclass of \texttt{BankAccount} to introduce minimum balance tracking.
\item Defining \texttt{ChoiceQuestion} as a subclass of \texttt{Question} to manage selectable options and numbering.
\item Creating subclasses \texttt{SedanCar} and \texttt{CompactCar} when they differ solely in numerical fuel capacity values.
\item Defining \texttt{Manager} as a subclass of \texttt{Employee} to incorporate department bonus calculations.
\end{enumerate}
\end{enumerate}
\vspace{-1mm}
\section*{Section B: Output Tracing, Inheritance Mechanics \& Diagnostics \hfill \normalsize\normalfont(25 points)}
\textit{Answer Q11--Q15. Each question is worth 5 points. Suggested time: 15 minutes.}
\subsection*{Q11. Exact Output of Multi-Class Inheritance \& Dynamic Dispatch Tracing (5 points)}
Predict the \textbf{exact console output} of the following Python program, including brackets, spacing, and capitalization.
\begin{lstlisting}
class Device:
    def __init__(self, brand, power): self._brand, self._power = brand, power
    def specs(self): return f"{self._brand}: {self._power}W"
    def status(self): return f"Device[{self.specs()}]"

class SmartDevice(Device):
    def __init__(self, brand, power, ip):
        super().__init__(brand, power); self._ip = ip
    def specs(self): return f"{super().specs()}, IP={self._ip}"

d1 = Device("Generic", 50)
d2 = SmartDevice("Echo", 15, "192.168.1.10")
print(d1.status())
print(d2.status())
print(f"{isinstance(d2, Device)} | {isinstance(d1, SmartDevice)}")
\end{lstlisting}
\begin{answerbox}[height=2.5cm]{Output Box for Q11}\end{answerbox}
\newpage

% =========================================================================
% PAGE 3: Q12-Q15
% =========================================================================
\subsection*{Q12. Evaluate Class Hierarchy, Type Introspection \& Inheritance Expressions (5 points)}
Given \texttt{class A: pass; class B(A): pass; class C(B): pass; b = B(); c = C()}. Write evaluated literals or booleans:
\begin{enumerate}[label=\alph*),leftmargin=2em,itemsep=1pt]
\item Evaluated boolean result of \lstinline|isinstance(c, A)|: \hrulefill
\item Evaluated boolean result of \lstinline|issubclass(A, B)|: \hrulefill
\item Evaluated boolean result of \lstinline|type(b) is A|: \hrulefill
\item Evaluated boolean result of \lstinline|issubclass(C, (A, list))|: \hrulefill
\item Evaluated boolean result of \lstinline|isinstance(b, (str, C))|: \hrulefill
\end{enumerate}
\vspace{-2mm}
\subsection*{Q13. Diagnose and Correct Three Object-Oriented Inheritance Defects (5 points)}
For each defect, identify the underlying bug and provide the corrected Python code snippet.\par
\vspace{0.3mm}
\renewcommand{\arraystretch}{1.05}\small
\begin{tabularx}{\linewidth}{|p{4.8cm}|X|X|}\hline
\textbf{Erroneous Code} & \textbf{Problem / Root Cause} & \textbf{Correction} \\ \hline
{\ttfamily\scriptsize class Animal:}\newline{\ttfamily\scriptsize \ \ def \_\_init\_\_(self, name):}\newline{\ttfamily\scriptsize \ \ \ \ self.\_name = name}\newline{\ttfamily\scriptsize class Dog(Animal):}\newline{\ttfamily\scriptsize \ \ def \_\_init\_\_(self, name, breed):}\newline{\ttfamily\scriptsize \ \ \ \ self.\_breed = breed}\newline{\itshape Goal: initialize name \& breed} & & \\[3.0mm]\hline
{\ttfamily\scriptsize class Employee:}\newline{\ttfamily\scriptsize \ \ def get\_salary(self): return 5000.0}\newline{\ttfamily\scriptsize class Manager(Employee):}\newline{\ttfamily\scriptsize \ \ def get\_salary(self):}\newline{\ttfamily\scriptsize \ \ \ \ return self.get\_salary() + 1000.0}\newline{\itshape Goal: add \$1000 bonus to base} & & \\[3.0mm]\hline
{\ttfamily\scriptsize class Base:}\newline{\ttfamily\scriptsize \ \ def \_\_init\_\_(self): self.\_\_token = 99}\newline{\ttfamily\scriptsize class Sub(Base):}\newline{\ttfamily\scriptsize \ \ def show(self): return self.\_\_token}\newline{\itshape Goal: read token in subclass} & & \\[3.0mm]\hline
\end{tabularx}\normalsize
\vspace{-2mm}
\subsection*{Q14. Savings Account Minimum Balance State Tracing with Overriding (5 points)}
Predict exact console output of the code below and show \texttt{\_b} and \texttt{\_m} after each step.
\begin{lstlisting}
class Account:
    def __init__(self, b): self._b = b
    def withdraw(self, a): self._b -= a
class Savings(Account):
    def __init__(self, b, r=0.1): super().__init__(b); self._r, self._m = r, b
    def withdraw(self, a): super().withdraw(a); self._m = min(self._m, self._b)
    def interest(self): return self._m * self._r

s = Savings(1000); s.withdraw(300); s.withdraw(400)
print(f"Interest: {s.interest():.2f} | Balance: {s._b:.2f}")
\end{lstlisting}
\begin{answerbox}[height=1.5cm]{Output and Trace Table for Q14}\end{answerbox}
\vspace{-2mm}
\subsection*{Q15. Dynamic Method Dispatch vs Type Checking in Polymorphic Fleets (5 points)}
Analyze the code: (1) Output of the print statement; (2) Why polymorphism eliminates manual \lstinline|isinstance()| checks.
\begin{lstlisting}
class Car:
    def toll(self): return 5.0
class Bus(Car):
    def toll(self): return 12.0
class EV(Car):
    def toll(self): return super().toll() * 0.5   # 50% clean fuel discount
print([v.toll() for v in [Car(), Bus(), EV()]])
\end{lstlisting}
\begin{answerbox}[height=1.4cm]{Answers and Technical Explanation for Q15}\end{answerbox}
\newpage

% =========================================================================
% PAGE 4: SECTION C (DESIGN & CLASS DECOMPOSITION)
% =========================================================================
\section*{Section C: Object-Oriented Design \& Hierarchical Class Decomposition \hfill \normalsize\normalfont(25 points)}
\textit{Answer all four tasks. Suggested time: 15 minutes.}
\begin{tcolorbox}[colback=boxbg,colframe=secondaryteal,arc=1.5mm,boxrule=0.6pt,top=1.5mm,bottom=1.5mm]
\textbf{Problem Specification: Automated Online Assessment \& Examination Engine (\texttt{Question} Hierarchy)}\\[0.5mm]
\small
An academic e-learning platform deploys an extensible hierarchy of examination questions (Horstmann Section 10.1--10.3).
\begin{itemize}
\item \textbf{Base Class \texttt{Question}:} Maintains a question prompt \texttt{\_text} (str) and a correct answer \texttt{\_answer} (str).
  \begin{itemize}
  \item \texttt{setText(questionText)}: Assigns the question prompt string.
  \item \texttt{setAnswer(correctResponse)}: Sets expected answer string (stripped of leading/trailing whitespace).
  \item \texttt{checkAnswer(response)}: Case-insensitive check comparing \texttt{response.strip().lower()} to \texttt{self.\_answer.lower()}. Returns boolean.
  \item \texttt{display()}: Prints the question text \texttt{self.\_text} directly to console.
  \end{itemize}
\item \textbf{Subclass \texttt{ChoiceQuestion(Question)}:} Extends \texttt{Question} to manage multiple-choice options.
  \begin{itemize}
  \item Initializes inherited state via \texttt{super().\_\_init\_\_()} and creates an empty list \texttt{\_choices = []}.
  \item \texttt{addChoice(choice, correct=False)}: Appends \texttt{choice} to \texttt{\_choices}. If \texttt{correct} is \texttt{True}, registers option's 1-based index (e.g., \texttt{"1"}, \texttt{"2"}) as the correct answer via \texttt{self.setAnswer(str(len(self.\_choices)))}.
  \item Overrides \texttt{display()}: Invokes \texttt{super().display()}, then prints each choice prefixed by its 1-based number: \texttt{"1: choice"}, \texttt{"2: choice"}, etc.
  \end{itemize}
\item \textbf{Subclass \texttt{NumericQuestion(Question)}:} Extends \texttt{Question} to accept floating-point numbers within tolerance $\varepsilon$ (e.g., $0.01$). Overrides \texttt{setAnswer(expected\_val)} to store float value, and overrides \texttt{checkAnswer(response)} to parse float safely and return $\text{True}$ iff $|\text{float}(\text{response}) - \text{expected}| \le \varepsilon$.
\end{itemize}
\end{tcolorbox}
\vspace{0.5mm}
\textbf{Task 1: Specify Public Methods, Parameter Types, Return Types, and Inheritance Contracts (6 points)}
Specify the method signatures and behavioral contracts for \texttt{ChoiceQuestion}:
\begin{itemize}[leftmargin=1.8em,itemsep=1.5pt]
\item \texttt{\_\_init\_\_(self)} $\rightarrow$ Return Type: \rule{2.5cm}{0.4pt} \quad Base delegation contract: \dotfill
\item \texttt{addChoice(self, choice: str, correct: bool = False)} $\rightarrow$ Return Type: \rule{2.0cm}{0.4pt} \quad Invariant: \dotfill
\item \texttt{display(self)} $\rightarrow$ Return Type: \rule{2.5cm}{0.4pt} \quad Console output contract: \dotfill
\item Does \texttt{ChoiceQuestion} need to override \texttt{checkAnswer(response)}? Explain: \dotfill
\end{itemize}

\textbf{Task 2: Implement the Complete Subclass \texttt{ChoiceQuestion} (7 points)}
Write a modular Python class definition for \texttt{ChoiceQuestion(Question)}. Implement \texttt{\_\_init\_\_}, \texttt{addChoice}, and \texttt{display}.
\begin{answerbox}[height=3.8cm]{Python Implementation for Task 2}\end{answerbox}

\textbf{Task 3: Hand-Trace a Polymorphic Quiz Execution Sequence (6 points)}
Trace the execution of a quiz containing: $Q_1$ (\texttt{Question}: "Who created Python?", ans: "Guido van Rossum"), $Q_2$ (\texttt{ChoiceQuestion}: "First release year?", Choices: 1985, 1991 [correct], 2000), and $Q_3$ (\texttt{NumericQuestion}: "Value of 1/3?", ans: 0.33, $\varepsilon=0.01$).
\begin{enumerate}[label=\arabic*.,leftmargin=1.8em,itemsep=1pt]
\item User inputs \texttt{"guido van rossum"} for $Q_1$ $\rightarrow$ Method dispatched: \rule{3.0cm}{0.4pt} \quad Result (\texttt{True}/\texttt{False}): \rule{1.5cm}{0.4pt}
\item User inputs \texttt{"2"} for $Q_2$ $\rightarrow$ Method dispatched: \rule{3.0cm}{0.4pt} \quad Result (\texttt{True}/\texttt{False}): \rule{1.5cm}{0.4pt}
\item User inputs \texttt{"0.333"} for $Q_3$ $\rightarrow$ Method dispatched: \rule{3.0cm}{0.4pt} \quad Result (\texttt{True}/\texttt{False}): \rule{1.5cm}{0.4pt}
\end{enumerate}
Total Quiz Score: \rule{2.5cm}{0.4pt} \quad Did any question require manual \texttt{if/elif type} branching? \rule{2.0cm}{0.4pt}

\vspace{0.5mm}
\textbf{Task 4: Technical Analysis \& Object-Oriented Principles (6 points)}
\begin{enumerate}[label=\alph*),leftmargin=1.8em,itemsep=1.5pt]
\item In the polymorphic loop \lstinline|for q in quiz: q.display()|, explain how the Python runtime determines whether to execute \texttt{Question.display()} or \texttt{ChoiceQuestion.display()}.\\[-1pt]\rule{\linewidth}{0.4pt}
\item Why is inheriting from \texttt{Question} superior to creating an independent \texttt{ChoiceQuestion} class from scratch? Discuss code reuse and interface polymorphism.\\[-1pt]\rule{\linewidth}{0.4pt}
\end{enumerate}
\newpage

% =========================================================================
% PAGE 5: SECTION D, PROBLEM 1
% =========================================================================
\section*{Section D: Applied Programming in Python \hfill \normalsize\normalfont(30 points)}
\textit{Write complete, modular Python 3 programs. Use clean encapsulation, explicit superclass delegation with \texttt{super()}, proper parameter validation, and a \texttt{main()} demonstration driver. Standard library only. Suggested time: 18 minutes.}

\subsection*{Problem 1: Corporate Personnel \& Tiered Payroll Hierarchy (15 points)}
Implement an inheritance hierarchy modeling corporate compensation across standard employees, managers, and executives.
\textbf{Requirements:}
\begin{enumerate}[label=\arabic*.,leftmargin=2em,itemsep=1.2pt]
\item Base class \texttt{Employee}:
  \begin{itemize}
  \item \texttt{\_\_init\_\_(self, emp\_id, name, base\_salary)}: Validate that \texttt{emp\_id} is a non-empty string and \texttt{base\_salary >= 0.0} (raise \texttt{ValueError} if invalid). Store as \texttt{\_emp\_id}, \texttt{\_name}, \texttt{\_base\_salary} (float).
  \item Accessors: \texttt{get\_salary(self)} returns \texttt{\_base\_salary}; \texttt{get\_bonus(self)} returns default \texttt{0.0}.
  \item \texttt{total\_compensation(self)}: Returns \lstinline|self.get_salary() + self.get_bonus()|.
  \item \texttt{\_\_str\_\_(self)}: Returns \lstinline|f"Employee[#{self._emp_id}, {self._name}]: Total = ${self.total_compensation():.2f}"|.
  \end{itemize}
\item Subclass \texttt{Manager(Employee)}:
  \begin{itemize}
  \item \texttt{\_\_init\_\_(self, emp\_id, name, base\_salary, bonus=0.0)}: Explicitly call \texttt{super().\_\_init\_\_} for base attributes. Validate \texttt{bonus >= 0.0} and store in \texttt{\_bonus} (float).
  \item Mutator: \texttt{set\_bonus(self, bonus)} updates \texttt{\_bonus} (validating $\ge 0.0$).
  \item Overrides \texttt{get\_bonus(self)}: Returns \texttt{self.\_bonus}.
  \item Overrides \texttt{\_\_str\_\_(self)}: Returns \lstinline|f"Manager[#{self._emp_id}, {self._name}]: Salary = ${self.get_salary():.2f}, Bonus = ${self.get_bonus():.2f}, Total = ${self.total_compensation():.2f}"|.
  \end{itemize}
\item Subclass \texttt{Executive(Manager)}:
  \begin{itemize}
  \item \texttt{\_\_init\_\_(self, emp\_id, name, base\_salary, bonus=0.0, shares=0)}: Call \texttt{super().\_\_init\_\_}. Validate \texttt{shares >= 0} and store in \texttt{\_shares} (int).
  \item Mutator: \texttt{grant\_shares(self, count)} adds positive integer \texttt{count} to \texttt{\_shares}.
  \item Overrides \texttt{get\_bonus(self)}: Returns \lstinline|super().get_bonus() + self._shares * 5.0| (adds \$5.00 dividend per stock share).
  \item Overrides \texttt{\_\_str\_\_(self)}: Returns \lstinline|f"Executive[#{self._emp_id}, {self._name}]: Shares = {self._shares}, Total = ${self.total_compensation():.2f}"|.
  \end{itemize}
\item \texttt{main()}: Instantiate an \texttt{Employee} (\#E101, "Alice Smith", \$4500.00), a \texttt{Manager} (\#M202, "Bob Johnson", \$7000.00, bonus \$1500.00), and an \texttt{Executive} (\#X303, "Carol Danvers", \$12000.00, bonus \$4000.00, 500 shares). Store them in a list \texttt{staff}, iterate polymorphically to print each profile and compute total monthly payroll expenditure (\$31500.00). Catch a \texttt{ValueError} on negative salary.
\end{enumerate}
\begin{answerbox}[height=8.5cm]{Python Solution for Problem 1}\end{answerbox}
\newpage

% =========================================================================
% PAGE 6: SECTION D, PROBLEM 2
% =========================================================================
\subsection*{Problem 2: Polymorphic Bank Account Fee \& Interest Hierarchy (15 points)}
Implement a polymorphic banking hierarchy modeling specialized fee-based checking and interest-bearing savings accounts, directly applying Horstmann Section 10.4.
\textbf{Requirements:}
\begin{enumerate}[label=\arabic*.,leftmargin=2em,itemsep=1.0pt]
\item Base class \texttt{BankAccount}:
  \begin{itemize}[itemsep=0.5pt]
  \item \texttt{\_\_init\_\_(self, acc\_num, balance=0.0)}: Validate \texttt{balance >= 0.0}; store \texttt{\_acc\_num} (str) and \texttt{\_balance} (float).
  \item \texttt{deposit(self, amount)}: Validate \texttt{amount > 0.0}; add to \texttt{\_balance}.
  \item \texttt{withdraw(self, amount)}: Validate \texttt{amount > 0.0}. If \texttt{amount <= self.\_balance}, deduct and return \texttt{True}; else return \texttt{False}.
  \item \texttt{month\_end(self)}: Template method for end-of-month processing. Base class body executes \texttt{pass}.
  \item \texttt{get\_balance(self)}: Returns current balance. \texttt{\_\_str\_\_(self)}: Returns formatted string.
  \end{itemize}
\item Subclass \texttt{SavingsAccount(BankAccount)}:
  \begin{itemize}[itemsep=0.5pt]
  \item \texttt{\_\_init\_\_(self, acc\_num, balance=0.0, rate=6.0)}: Delegate via \texttt{super()}; store \texttt{\_rate} (float); initialize \texttt{\_min\_balance = self.\_balance}.
  \item Overrides \texttt{withdraw(self, amount)}: Call \lstinline|super().withdraw(amount)|. If successful and balance falls below minimum, update \texttt{\_min\_balance = self.\_balance}. Return success.
  \item Overrides \texttt{month\_end(self)}: Calculate monthly interest: $\text{interest} = \text{\_min\_balance} \times (\text{rate}/12) / 100$. Deposit interest, reset $\text{\_min\_balance} = \text{\_balance}$, and return \texttt{interest}.
  \end{itemize}
\item Subclass \texttt{CheckingAccount(BankAccount)}:
  \begin{itemize}[itemsep=0.5pt]
  \item \texttt{\_\_init\_\_(self, acc\_num, balance=0.0, free=3, fee=1.50)}: Delegate via \texttt{super()}; store \texttt{\_free} (int), \texttt{\_fee} (float), initialize \texttt{\_trans = 0}.
  \item Overrides \texttt{deposit(self, amount)} and \texttt{withdraw(self, amount)}: Delegate to \texttt{super()} and increment \texttt{\_trans} upon success.
  \item Overrides \texttt{month\_end(self)}: Deduct excess fees $\max(0, \text{\_trans} - \text{\_free}) \times \text{\_fee}$ from balance. Reset $\text{\_trans} = 0$, and return fee charged.
  \end{itemize}
\item \texttt{main()}: Create \texttt{SavingsAccount("SA-101", 1200.0, 6.0)} and \texttt{CheckingAccount("CA-202", 500.0, free=2, fee=2.00)}. Perform deposit and withdrawal operations. In a polymorphic loop \lstinline|for acc in [sa, ca]: acc.month_end()|, execute monthly close and verify final balances (\$1005.00 and \$496.00).
\end{enumerate}
\begin{answerbox}[height=8.5cm]{Python Solution for Problem 2}\end{answerbox}
\vfill
\begin{center}\rule{0.5\textwidth}{0.4pt}\\[1mm]
\textbf{\color{primaryblue}--- END OF EXAM PAPER ---}
\end{center}
\end{document}
